% paper/abstract_completeness.tex
% Option A — completeness-first narrative. Use this when sim results
% do NOT show a clear task_success win for GRACE. Honest, low-risk.
% ~200 words.

\begin{abstract}
Large language models are now the default cognitive core of embodied
household agents, yet a closer look at the plans they produce reveals
a hidden trade-off: high task-completion rates are typically sustained
by execution-time retry loops that paper over plans that fail
symbolic precondition checks. We propose \textbf{GRACE} (Grounded
Repair-and-Critique Embodied planner), a single-LLM planner that
integrates three lightweight components: (i) a rule-based symbolic
precondition verifier ($\sim$250 lines of Python, zero tokens,
$O(|\pi|)$ time) that gates each sub-plan; (ii) hierarchical
decomposition with sub-goal--localised repair that regenerates only
the suffix of the failed sub-goal; and (iii) a hybrid memory store
seeded from curated ground-truth plans and grown by successful
runtime episodes. Evaluated on the 38-task AI2-THOR household
benchmark across three open-weight Ollama models, GRACE achieves the
highest \emph{goal completeness} of any of the eight pyplanner
baselines on all three models tested ($+0.10$--$0.13$ over the best
baseline), with a controlled increase in cost. Ablations isolate the
three components' contributions. We discuss the cost--completeness
trade-off honestly: GRACE consumes more tokens than one-shot
baselines but is the only method that consistently surfaces every
expected interactable in long-horizon plans. Code and benchmark
scripts are released as a self-contained artifact.
\end{abstract}
